\begin{lemma}[Absorbing vertices; Scheide Proposition 3.5]
If $(T,e_0,\varphi)$ is globally maximum, then
\[
  V(T)\cup A(T,e_0,\varphi)
\]
is elementary with respect to $\varphi$.
\end{lemma}
\begin{proof}
Put $X=V(T)$ and $A=A(T,e_0,\varphi)$.  The statement-level package
\noderef{ScheideGlobalMaximumOrder} says that $X$ is elementary and closed.
It remains to compare an absorbing vertex with $X$, and then two distinct
absorbing vertices.

Suppose first that $v\in A$, that $\delta$ is missing at $v$, and that
$\delta$ is also missing at some $x\in X$.  By
\noderef{ScheideGlobalMaximumFourFree} there is a free color
$\gamma\ne\delta$.  Elementariness gives a unique $u\in X$ missing
$\gamma$.  Since both $\gamma$ and $\delta$ are missing on the closed set
$X$, closedness forbids every boundary edge identity of either color.  A
walk in the two-color incidence graph from the outside vertex-node $v$ to
a vertex-node of $X$ would have a first edge-node crossing the boundary of
$X$; that edge identity would have color $\gamma$ or $\delta$, a
contradiction.  Hence the $(\gamma,\delta)$-component through $v$ is
disjoint from $X$.  This contradicts the defining requirement in
\noderef{AbsorbingVertex} that it contain $u$.  Thus no color missing at a
vertex of $A$ is missing on $X$.

Now let $v,w\in A$ be distinct and suppose that $\delta$ is missing at
both.  Again choose a free $\gamma\ne\delta$, and let $u\in X$ be the
unique vertex missing $\gamma$.  The first paragraph shows that $\delta$
is present at $u$.  Properness, supplied by the $k$-triple clause of
\noderef{GlobalMaximumTriple}, gives degree one at each of the three
vertex-nodes $u,v,w$ in their common $(\gamma,\delta)$-component: $u$
misses $\gamma$ and has its unique $\delta$-edge, whereas $v,w$ miss
$\delta$ and have their unique $\gamma$-edges.  The absorbing property
puts all three nodes in that component.  But
\noderef{TwoColorComponentSpanningPath}, together with the two-color
degree bound used there, says the finite component is a path; a path has
only two degree-one vertices.  This is impossible.  Missing-color sets of
distinct vertices of $X\cup A$ are therefore pairwise disjoint, which is
exactly elementariness as defined in
\noderef{ColoredVertexSetElementary}.
\end{proof}
